\documentclass[10pt,a4paper]{article}

% 1. Kodowanie i języki
\usepackage[T1]{fontenc}
\usepackage[utf8]{inputenc}
\usepackage[polish]{babel}

% 2. Font o stałej szerokości (Monospaced) dla CAŁEGO dokumentu
\usepackage{inconsolata}
\renewcommand{\familydefault}{\ttdefault}

% 3. Marginesy - maksymalne upakowanie tekstu (0.5 cala)
\usepackage[margin=0.5in]{geometry}

% 4. Interlinia i formatowanie akapitów
\linespread{0.9}
\setlength{\parindent}{0pt}
\setlength{\parskip}{4pt}

% 5. Kolory do wizualnego rozróżnienia ról
\usepackage{xcolor}
\definecolor{rulecolor}{gray}{0.7}

% 6. Proste makra do oznaczania nadawcy (poprawione definicje)
\newcommand{\UserTurn}{\textbf{[USER]}\par\medskip}
\newcommand{\ModelTurn}{\textbf{[GEMINI]}\par\medskip}
\newcommand{\Separator}{\par\vspace{4pt}\noindent{\color{rulecolor}\rule{\linewidth}{0.4pt}}\par\vspace{6pt}}

\begin{document}
\UserTurn
Czesc Gemini. Chce przetestowac jak zapisujesz kontekst rozmowy w jsonie - no wiesz - rola 'model' i 'user'. Daj tutaj jakis tekst jako odpowiedz, niech bedzie reprezentatywny w miare, o jakims poziomie skomplikowania pod wzgledem zapisu (np  znaki specjalne, zagniezdzenia) ale niech to nie bedzie  jakis straszny corner case, worst case. Pozostawiam Tobie pole do wykazania sie.
\Separator
\ModelTurn
Cześć! Bardzo chętnie pomogę Ci w tym teście. 

Przygotowałem dla Ciebie tekst o umiarkowanym stopniu skomplikowania. Znajdziesz w nim polskie znaki diakrytyczne, podstawowe znaczniki formatowania (np. cudzysłowy, apostrofy), proste zagniezdzenia w postaci listy oraz fragment kodu z nawiasami klamrowymi i ukośnikami, które często bywają wyzwaniem przy parsowaniu JSON-a.

Oto treść wiadomości:

***

Analiza zapytania wykazała następujące parametry wejściowe dla modułu "Główny\_Procesor\_V2":

* Status: OK (kod: `200\_SUCCESS`)
* Ścieżka dostępu: `/api/v1/user/settings\#auth`
* Komunikat zwrotny: „Operacja zakończona pomyślnie. Zaktualizowano preferencje użytkownika (w tym: język, strefa czasowa oraz tokeny sesji).”

Warto zwrócić uwagę na sekcję konfiguracyjną, która zawiera zagnieżdżone obiekty:
\begin{verbatim}
{
  "system": {
    "tryb_debugowania": false,
    "timeout_ms": 5000,
    "dozwolone_domeny": ["*.example.com", "api.test-env.pl"]
  }
}
\end{verbatim}

Proszę o potwierdzenie odbioru danych oraz informację, czy powyższa struktura spełnia Twoje oczekiwania odnośnie testu kontekstu.

***

Jestem ciekaw, jak Twój skrypt poradzi sobie z zapisem tego do pliku JSON. Daj znać, jak poszło!
\Separator
\UserTurn
Zawartosc Twojej odpowiedzi od slowa 'Czesc' do 'jak poszlo' widze w trybie tekstowym. Zapewne wiesz ze w pliku json bedacym kontekstem tej naszej konwersacji jest to zapisane w taki sposob mniej wiecej:'role' model i 'text': Czesc... jak poszlo. Ten plik json potraktuje jako modelowe wymaganie dla parsera ktorego chce stworzyc. Dobrze sie zrozumielismy, o to Ci chodzilo?
\Separator
\ModelTurn
Tak, doskonale się zrozumieliśmy! 

Dokładnie o to chodzi – moja poprzednia odpowiedź (wraz z całym jej "bogatym" w znaki specjalne, cudzysłowy, znaczniki Markdown i wcięciami wnętrzem) trafiła do pola `text` w obiekcie o `role: "model"` w strukturze JSON zarządzającej naszą konwersacją.

Dla Twojego parsera oznacza to, że musi on sobie radzić z wyzwaniami takimi jak:
1. **Wielolinijkowość** – znaki nowej linii (`\textbackslash\{\}n`) wewnątrz wartości stringa.
2. **Cudzysłowy i apostrofy** – które wewnątrz tekstu JSON muszą być odpowiednio esejowane (np. `\textbackslash\{\}"`), aby nie "zamknąć" przedwcześnie wartości klucza `text`.
3. **Znaki specjalne i formatowanie** – polskie litery (UTF-8), bloki kodu z nawiasami klamrowymi `\{\}` oraz ukośniki `/`.

Jeśli Twój parser poprawnie wyciągnie tę treść i z powrotem złożone drzewo JSON nie wywoła błędów składniowych (np. `JSON.parse()` w JavaScript czy `json.loads()` w Pythonie), to znaczy, że konstrukcja jest jak najbardziej poprawna.

Powodzenia w budowaniu parsera! Daj znać, czy napotkasz jakieś pułapki po drodze.
\Separator
\UserTurn
Dziekuje. Dobra robota.
\Separator
\ModelTurn
Cieszę się, że mogłem pomóc! 

Gdybyś potrzebował jeszcze jakichś danych testowych do parsera (np. z jeszcze bardziej złośliwymi znakami, jak emoji, znaki kontrolne czy zagnieżdżone JSON-y wewnątrz JSON-a) – śmiało wołaj. 

Powodzenia w kodowaniu i miłego dnia!
\Separator
\end{document}
